\usepackage[T1]{fontenc}
\usepackage[utf8]{inputenc}
\usepackage[brazil]{babel}
\usepackage{lmodern}
\usepackage{microtype}

\usepackage{amsmath,amssymb,mathtools}
\usepackage{siunitx}
\sisetup{
  locale = DE,
  per-mode = symbol,
  separate-uncertainty = true
}

\usepackage{booktabs}
\usepackage{graphicx}
\usepackage{xcolor}
\usepackage{listings}
\usepackage{caption}
\usepackage{subcaption}

\usepackage[
  colorlinks=true,
  linkcolor=blue!50!black,
  citecolor=green!40!black,
  urlcolor=blue!60!black
]{hyperref}
\usepackage[nameinlink,noabbrev,brazilian]{cleveref}

\graphicspath{{figuras/}}

\newcommand{\vetor}[1]{\boldsymbol{#1}}
\newcommand{\dd}{\mathop{}\!\mathrm{d}}
\newcommand{\e}{\mathrm{e}}

\definecolor{corcodigo}{RGB}{30,70,105}
\definecolor{fundocodigo}{RGB}{247,249,251}

\lstdefinestyle{pythonlivro}{
  language=Python,
  basicstyle=\ttfamily\small,
  keywordstyle=\color{corcodigo}\bfseries,
  commentstyle=\color{green!40!black},
  stringstyle=\color{orange!60!black},
  backgroundcolor=\color{fundocodigo},
  frame=single,
  rulecolor=\color{black!20},
  numbers=left,
  numberstyle=\tiny\color{black!55},
  breaklines=true,
  showstringspaces=false,
  columns=fullflexible,
  keepspaces=true
}

\lstset{style=pythonlivro}

